\documentclass[a4paper,10pt, oneside]{book}
\usepackage[utf8]{inputenc}
\usepackage[spanish]{babel}
\usepackage[style=ieee,backend=bibtex]{biblatex}
\usepackage{graphicx}
\usepackage{draftwatermark}
\usepackage{lineno}
\usepackage[table]{xcolor}
%\usepackage{showframe}
\usepackage[a4paper]{geometry}
\usepackage{caption}
\usepackage{pgfgantt}
\usepackage{lscape} 


\bibliography{anteproyecto}

\SetWatermarkText{}
\SetWatermarkScale{5}
\SetWatermarkLightness{0.95}

\renewcommand{\labelenumii}{\theenumii}
\renewcommand{\theenumii}{\theenumi.\arabic{enumii}.}


	
\begin{document}
\begin{titlepage}
	\centering
	\includegraphics[width=0.25\textwidth]{Universidad_del_Litoral}\par\vspace{1cm}
	{\scshape\LARGE Universidad Nacional del Litoral \par}	
	%\vspace{1cm}
	{\scshape\Large Proyecto Final de Carrera\par}
	\vspace{1.5cm}
	{\huge\bfseries Diseño de un sistema de detección de anomalías en redes de computadoras.\par}
	\vspace{6cm}
	realizado por \par
	{\Large\itshape Leandro Pineda\par}
	\vspace{1cm}
	dirigido por Ing. Miguel Angel Robledo\par
	codirigido por Ing. Gabriel Filippa
	

	\vfill
	
	%Repositorio git del documento (solo versión borrador): https://github.com/leandropineda/anteproyecto
	
	%\vfill
	
	% Bottom of the page
	\large Santa Fe\par
	{Septiembre de 2016 \par}
	
\end{titlepage}

\modulolinenumbers[5]
%\linenumbers

\section*{Resumen}

En la actualidad las redes informáticas están presentes en las operaciones diarias de prácticamente cualquier empresa. Además hacemos uso de Internet a diario ya que cada vez se ofrecen más servicios en linea. Por esto, las herramientas que facilitan la tarea de administración y diagnóstico de redes son de gran utilidad para ofrecer un servicios de mejor calidad, y aprovechar la utilización de recursos existentes.
El presente documento describe un proyecto para el desarrollo de un sistema de monitoreo y visualización en tiempo real del tráfico de red que atraviesa una interfaz de red, que puede ser puesto en funcionamiento con un bajo impacto en la infraestructura productiva de cualquier empresa, con un bajo impacto y a un bajo costo.

\paragraph{Palabras claves} \textit{tiempo real}, \textit{redes informáticas}, \textit{Apache Storm}, \textit{monitoreo}, \textit{visualización}

\newpage
\section*{Justificación}
Las redes informáticas son esenciales para que las empresas o instituciones públicas puedan realizar sus tareas diarias. A medida que estas entidades crecen también lo hace el volumen de datos que generan, demandando así una infraestructura de red con mayor tamaño. Debido a esto, mantener una buena performance en el funcionamiento general de un sistema que se encuentra constantemente en crecimiento se convierte en una tarea desafiante.

Algunas de las herramientas disponibles para monitoreo de redes brindan reportes basados en estadísticas. La empresa Cisco\textsuperscript{\textregistered} introduce una característica en algunos modelos de routers y switches llamada NetFlow, la cual permite recolectar información sobre el tráfico de red que atraviesa las distintas interfaces del dispositivo. Sin embargo, estar limitado exclusivamente a modelos particulares de routers y switches significa una desventaja pues en algunas ocasiones no se dispone del hardware mencionado, lo que implica una inversión en infraestructura considerable. Otros fabricantes introducen características similares a NetFlow, pero con nombres diferentes: Traffic Flow de MikroTik, NetStream de HP, entre otros.
Otra alternativa para monitorear el tráfico que pasa a través de una interfaz de red es el analizador de paquetes \textit{tcpdump} (disponible en sistemas GNU/Linux), el cual permite observar una descripción en texto plano del contenido de los paquetes que atraviesan la interfaz.

Disponer de la información más actualizada posible de los eventos que han ocurrido en la red facilita la tarea de administrar y diagnosticar fallas. Es por esto que el monitoreo en tiempo real es un herramienta de gran utilidad pues permite contar con información de una serie de métricas muy descriptivas al instante. El proceso de toma de decisiones a la hora de invertir en infraestructura puede realizarse de forma adecuada si se incluye la información que puede extraerse observando la utilización de la infraestructura actual. Por ejemplo, una empresa podría decidir si invertir en una conexión a Internet de mayor velocidad o limitar el ancho de banda disponible para ciertos servicios que hacen uso intensivo de las redes. El monitoreo y la visualización en tiempo real puede hacer también que el proceso de detección de problemas en una red sea rápido y sencillo.

Se propone entonces diseñar un sistema capaz de registrar los eventos ocurridos en capa de transporte y mostrar información acerca del tráfico de red que está siendo transmitido por una red. El mismo desarrollará utilizando técnicas llamadas \textit{Real-Time Stream Processing} para analizar la salida de texto plano del analizador de paquetes \textit{tcpdump} y producir en tiempo real informes acerca de los parámetros de performance más importantes\footnote{IP Performance Metrics Working Group https://datatracker.ietf.org/wg/ippm/documents/ (IPPM)}
. El uso de estas técnicas permite diseñar un sistema escalable el cual pueda adaptarse a redes que están en crecimiento constante.

Sin duda alguna, la cantidad de dispositivos conectados a Internet crece día a día y con ellos lo hace la demanda de la infraestructura de red necesaria para brindar un servicio de calidad. Contar con información del uso de las redes es fundamental para hacer estimaciones para inversiones a futuro, y permite determinar si los recursos disponibles están siendo utilizados de manera eficiente.

Finalmente, contar con esta información es una gran ayuda para mejorar la seguridad de las redes. Muchos de los ataques a redes informáticas son realizados por usuarios fuera de la red que no tienen acceso físico a la misma, pero muchos de los ataques más peligrosos son llevados a cabo por usuarios internos. Tener este tipo de información sobre el tráfico de la red puede ser una herramienta de enorme valor para identificar estos sucesos.

\section*{Objetivos}
\subsection*{Generales}
Diseñar un sistema que permita monitorear en tiempo real parámetros de performance del tráfico que atraviesa una interfaz de red.
\subsection*{Específicos}
\begin{itemize}

\item Identificar un conjunto de tecnologías y herramientas adecuadas para el diseño del sistema.
\item Diseñar un sistema de monitoreo de redes que funcione en \textit{commodity hardware}.
\item Construir una herramienta de software que registre los eventos ocurridos en la red (capa de transporte).
\item Elaborar un documento que describa la arquitectura del sistema.
\end{itemize}

\section*{Alcances}

\subsection*{Funcionales}
\begin{itemize}
	\item El sistema permitirá visualizar en tiempo real los parámetros de performance de la red. Se podrá monitorear el tráfico a nivel de \textit{host} (capa de red) y a nivel de servicios (capa de transporte)
	\item Los informes mostrarán información de la Capa de Transporte y la Capa de Red\footnote{Modelo TCP/IP}. \cite{rfc791}\cite{rfc793}\cite{rfc1180}
	\item El sistema proveerá al usuario una interfaz para la visualización de los informes. Estos contaran con los siguientes parámetros:
	\begin{itemize}
		\item Utilización de la red
		\item Paquetes por segundo transmitidos
		\item Cantidad de peticiones de conexión
		\item Cantidad de conexiones realizadas
		\item Cantidad de conexiones rechazadas y expiradas
		\item Cantidad de conexiones activas
		\item Tráfico entrante y saliente por host
		\item Tráfico entrante y saliente por proceso
	\end{itemize}
	\item El sistema almacenará información de las conexiones establecidas. Lo hará en archivos de texto plano y en formato semi-estructurado.
	\item El sistema no identificará patrones o generará alarmas.
\end{itemize}

\subsection*{No Funcionales}
\begin{itemize}
	\item El sistema será escalable. Las tecnologías seleccionadas permiten configurar un cluster de nodos de procesamiento en caso que el volumen de datos a procesar aumente.
	\item Se proveerá un manual de instalación y configuración.
\end{itemize}
 
\subsubsection{Exclusiones}
El proyecto no contempla la instalación del sistema en un ámbito de producción. El sistema a desarrollar solo analizará el tráfico que atraviesa una única interfaz de red, y no contempla el caso de conexiones múltiples con balance de carga, ni la utilización de múltiples nodos configurados en modo cluster. No se implementará ningún mecanismo de detección ataques o de patrones de comportamiento sospechoso.

\subsubsection{Supuestos}
Los datos necesarios para realizar pruebas pueden ser generados arbitrariamente.
El término \textit{tiempo real} refiere al denominado \textit{tiempo real blando}, es decir, no es necesario asegurar la ejecución de ciertas tareas o mostrar informes en el mismo instante en el que se generan los datos. La generación de los informes tendrá entonces demoras de cientos de milisegundos.

\subsubsection{Criterios de Aceptación}
Se considera que el proyecto está aceptado cuando cumple en un 90\% los requisitos funcionales, con un nivel mínimo de aceptación del 75\%. Los prototipos deben tener implementadas todas las funcionalidades planificadas.
\section*{Metodología}

La metodología de desarrollo del proyecto será incremental e iterativa. Incremental porque varios componentes y funcionalidades del sistema se desarrollarán en momentos diferentes y serán integradas cuando sean completadas. Iterativa pues se invertirán esfuerzos en revisar constantemente partes del sistema, tanto para mejorar la calidad externa como interna del software\cite{ISOIEC9126}.

Dado que los requerimientos de los interesados pueden cambiar en el transcurso de la ejecución del proyecto, se propone utilizar un enfoque de desarrollo ágil. El mismo contempla la posibilidad de priorización y selección en los alcances y en las prioridades de las diferentes funcionalidades, y fundamentalmente promueve la entrega continua de software y la inclusión de los interesados en el proceso de desarrollo.

\section*{Plan de Tareas}

El proyecto se divide en 5 incrementos. A continuación se da una breve descripción de los mismos:

\paragraph{Incremento 1: Investigación preliminar}\
La primer parte del proyecto consiste en determinar que conjunto de tecnologías serán utilizadas, y elaborará una descripción a alto nivel de las diferentes componentes del sistema. Además se realizará una investigación sobre las técnicas de modelado de comportamiento existentes con el fin de determinar cuales serán implementadas.
\paragraph{Incremento 2: Modelado de comportamiento} \
En esta etapa se implementarán las técnicas de modelado de comportamiento seleccionadas en el incremento anterior. Se utilizarán solo los datos provenientes de la capa de transporte y se desarrollará un software con una interfaz web sencilla donde se muestren los principales parámetros del modelado.
\paragraph{Incremento 3: Detección de anomalías} \
Se implementarán las funcionalidades de detección de anomalías y se agregará al software mecanismos de generación de alarmas. Además se mejorará la interfaz de usuario.
\paragraph{Incremento 4: Modelado de comportamiento de subsistemas} \
Con el fin de mejorar la identificación de anomalías se incorporará al modelo de comportamiento información de los diferentes subsistemas que componen la infraestructura de red.
\paragraph{Incremento 5: Pruebas exhaustivas y documentación}
Se completará el desarrollo del software y se redactará la documentación necesaria. Dado que la tasa de falsos negativos\footnote{La existencia de falsos negativos implica la existencia de ataques verdaderos que no son detectados por el IDS} debe ser mínima, se realizarán pruebas exhaustivas de funcionamiento. Una vez concluidas las pruebas, se elaborará un informe con los resultados obtenidos

\ \

Al finalizar la primera iteración de cada incremento se obtiene una herramienta de software con las funcionalidades descriptas y calidad de producto final, con excepción de la etapa de investigación preliminar donde se obtendrá un informe en soporte escrito o digital.



\begin{table}[htbp]
	\begin{center}	
		\begin{tabular}{|l|c|}
			\hline 
			Entregable & Fecha de entrega \\ \hline
			Informe de avance 1 & 04/10/2016 \\
			Informe de avance 2 & 08/11/2016 \\
			Informe de avance 3 & 16/12/2016 \\
			Informe de avance 4 & 14/02/2017 \\ \hline
		\end{tabular}
	\end{center}
	\caption{Fechas de entrega de informes de avance.}
	\label{table:informes}
\end{table}

\subsubsection*{Plan de Tareas}

La duración total del proyecto es de $466$ horas, con una dedicación de $20$ horas semanales. A continuación se detalla el plan de tareas.

\begin{enumerate}
	\setlength{\itemsep}{0pt}
	\setlength{\parskip}{0pt}
	\item \textbf{Investigación preliminar} (66hs)
	\begin{enumerate}
		\item Estudio comparativo de las tecnologías y métodos de modelado. (24hs)
		\item Diseño conceptual del sistema. (30hs)
		\item Documentación. (12hs)
	\end{enumerate}
	\item \textbf{Modelado de comportamiento} (84hs)
	\begin{enumerate}
		\item Instalación y configuración de la plataforma de desarrollo. (12hs)
		\item Implementación de funcionalidad de captura de tráfico de red. (24hs)
		\item Implementación de modelado de tráfico. (24hs)
		\item Implementación de interfaz web. (24hs)
	\end{enumerate}
	\item \textbf{Detección de anomalías} (98hs)
	\begin{enumerate}
		\item Implementación de funcionalidad de detección de anomalías. (30hs)
		\item Pruebas de implementación. (24hs)
		\item Implementación de módulo de alarmas. (24hs)
		\item Mejora de interfaz web. (20hs)
	\end{enumerate}
	\newpage
	\item \textbf{Modelado de comportamiento de subsistemas} (146hs)
	\begin{enumerate}
		\item Implementación de funcionalidad de captura de logs de subsistemas. (72hs)
		\begin{itemize}
			\item Servidores web (24hs).
			\item Gestores de base de datos (24hs).
			\item Firewalls (24hs).
		\end{itemize}
		\item Implementación de modelado de comportamiento de subsistemas. (30hs)
		\item Pruebas de implementación. (24hs)
		\item Implementación de funcionalidad de detección de anomalías y alarmas. (20hs)
	\end{enumerate}
	\item \textbf{Pruebas exhaustivas y documentación} (72hs)
	\begin{enumerate}
		\item Pruebas de detección. (20hs)
		\item Elaboración de informe de desempeño. (12hs)
		\item Elaboración de informe final. (40hs)
	\end{enumerate}
\end{enumerate}

\subsection*{Informes de avance}

Se presentarán 4 informes de avance en las fechas de finalización de cada etapa, detalladas en el Cuadro \ref{table:cronograma}. A continuación se detalla que información será incluida en cada informe:

\paragraph{Informe de avance 1} 
Contendrá los resultados obtenidos en los estudios comparativos de las tecnologías y las técnicas de modelado y detección y justificará la elección de las mismas. Además se incluirá una descripción general de la arquitectura del sistema.

\paragraph{Informe de avance 2}
Contendrá información sobre los criterios de selección de características para el modelado de tráfico de red. Se proveerá la guía de instalación y configuración de la plataforma. Además tendrá información sobre cambios realizados en los entregables anteriores y el desempeño del modelo implementado.

\paragraph{Informe de avance 3}
Contendrá información sobre los criterios de selección de características para el modelado del comportamiento de los subsistemas, y el desempeño del modelo implementado. Además se detallarán los cambios realizados en los entregables anteriores.

\paragraph{Informe de avance 4}
Se describirán las pruebas de integración del sistema. Además tendrá información sobre cambios realizados en los entregables anteriores y el desempeño general del sistema. También se incluirán los resultados de las pruebas de detección de la última iteración.

\begin{table}[htbp]
	\begin{center}	
		\begin{tabular}{|l|c|c|c|}
			\hline 
			Etapa & Inicio & Finalización & Duración \\ \hline
			Investigación preliminar & 01/09/2016 & 30/09/2016 & 4 semanas \\
			Modelado de comportamiento & 03/10/2016 & 04/11/2016 & 4 semanas \\
			Detección de anomalías & 07/11/2016 & 16/12/2016 & 5 semanas \\
			Modelado de comportamiento de subsistemas & 19/12/2016 & 18/02/2017 & 8 semanas \\
			Pruebas exhaustivas y documentación & 20/02/2017 & 25/03/2017 & 5 semanas \\ \hline
		\end{tabular}
	\end{center}
	\caption{Fechas estimativas de inicio y fin de actividades.}
	\label{table:cronograma}
\end{table}

\newpage

\subsection*{Diagrama de Gantt}

\begin{figure}[hb]
\begin{center}
\begin{ganttchart}[
	time slot format=isodate,
	x unit=.6mm,
	y unit title=.6cm,
	y unit chart=4mm,
	hgrid,
	link/.style={-latex}
	]{2016-09-01}{2017-03-29}
	\gantttitlecalendar{year, month} \\
	\ganttgroup{Tarea 1}{2016-09-01}{2016-09-30} \\
	\ganttbar[name=a11]{Actividad 1.1}{2016-09-01}{2016-09-08} \\
	\ganttbar[name=a12]{Actividad 1.2}{2016-09-10}{2016-09-30} \\
	\ganttbar[name=a13]{Actividad 1.3}{2016-09-10}{2016-09-30} \\
	\ganttlink{a11}{a12} 
	\ganttlink{a11}{a13} 
	\ganttgroup{Tarea 2}{2016-10-03}{2016-11-04} \\
	\ganttbar[name=a21]{Actividad 2.1}{2016-10-03}{2016-10-06} \\
	\ganttbar[name=a22]{Actividad 2.2}{2016-10-08}{2016-10-16} \\
	\ganttbar[name=a23]{Actividad 2.3}{2016-10-18}{2016-11-04} \\
	\ganttbar[name=a24]{Actividad 2.4}{2016-10-18}{2016-11-04} \\
	\ganttlink{a21}{a22}
	\ganttlink{a22}{a23}
	\ganttlink{a22}{a24} 
	\ganttgroup{Tarea 3}{2016-11-07}{2016-12-16} \\
	\ganttbar[name=a31]{Actividad 3.1}{2016-11-07}{2016-11-15} \\
	\ganttbar[name=a32]{Actividad 3.2}{2016-11-17}{2016-11-26} \\
	\ganttbar[name=a33]{Actividad 3.3}{2016-11-28}{2016-12-16} \\
	\ganttbar[name=a34]{Actividad 3.4}{2016-11-28}{2016-12-16} \\
	\ganttlink{a31}{a32}
	\ganttlink{a32}{a33}
	\ganttlink{a32}{a34}
	\ganttgroup{Tarea 4}{2016-12-19}{2017-02-18} \\
	\ganttbar[name=a41]{Actividad 4.1}{2016-12-19}{2017-01-17} \\
	\ganttbar[name=a42]{Actividad 4.2}{2017-01-19}{2017-02-10} \\
	\ganttbar[name=a43]{Actividad 4.3}{2017-01-19}{2017-02-10} \\
	\ganttbar[name=a44]{Actividad 4.4}{2017-02-12}{2017-02-18} \\
	\ganttlink{a41}{a42}
	\ganttlink{a41}{a43}
	\ganttlink{a42}{a44}
	\ganttlink{a43}{a44}
	\ganttgroup{Tarea 5}{2017-02-20}{2017-03-25} \\
	\ganttbar[name=a51]{Actividad 5.1}{2017-02-20}{2017-02-27} \\
	\ganttbar[name=a52]{Actividad 5.2}{2017-03-06}{2017-03-25} \\
	\ganttbar[name=a53]{Actividad 5.3}{2017-03-06}{2017-03-25}
	\ganttlink{a51}{a52}
	\ganttlink{a51}{a53}
\end{ganttchart}
\end{center}
\caption{Diagrama de Gantt del proyecto}
\end{figure}
%\end{landscape}

\section*{Riegos}

Para realizar el análisis cualitativo de los riesgos se asigna una probabilidad de ocurrencia y un impacto a cada riesgo. Luego se priorizan según su \textit{severidad}. Los riesgos identificados son:

\begin{center}
\begin{tabular}{|l||c|c|c|}
\hline
Riesgo & Imp. & \% ocur. & Sev. \\
\hline
No se pueden satisfacer las restricciones de performance.  & 5 & 3 & 15 \\
No se dispone del hardware necesario. & 4 & 3 & 12 \\
Los módulos del sistema no pueden ser integrados correctamente. & 4 & 2 & 8 \\
Las tecnologías seleccionadas no pueden ser integradas. & 4 & 2 & 8 \\
El sistema no puede ser probado en un entorno productivo. & 2 & 2 & 4 \\

\hline
\end{tabular}
\end{center}

\subsection*{Análisis Riesgos}
Del análisis de la matriz de probabilidad e impacto se deciden las estrategias.

\begin{center}
\renewcommand{\arraystretch}{1.2}
\begin{tabular}{|c||c|c|c|c|c|}

	\hline
	& 1 & 2 & 3 & 4 & 5 \\ \hline
	\hline
	1 & \cellcolor[gray]{0.8} & \cellcolor[gray]{0.8} & \cellcolor[gray]{0.8} & \cellcolor[gray]{0.8} & \cellcolor[gray]{0.6} \\ \hline
	2 & \cellcolor[gray]{0.8} & \cellcolor[gray]{0.8} & \cellcolor[gray]{0.6} & \cellcolor[gray]{0.6} & \cellcolor[gray]{0.6} \\ \hline
	3 & \cellcolor[gray]{0.8} & \cellcolor[gray]{0.6} & \cellcolor[gray]{0.6} & \cellcolor[gray]{0.6} & \cellcolor[gray]{0.4} \\ \hline
	4 & \cellcolor[gray]{0.8} & \cellcolor[gray]{0.6} & \cellcolor[gray]{0.6} & \cellcolor[gray]{0.4} & \cellcolor[gray]{0.4} \\ \hline
	5 & \cellcolor[gray]{0.6} & \cellcolor[gray]{0.6} & \cellcolor[gray]{0.4} & \cellcolor[gray]{0.4} & \cellcolor[gray]{0.4} \\ \hline
\end{tabular}
\qquad \qquad
\begin{tabular}{|c|l|}
	\hline Severidad & Estrategia \\ \hline
	menor que 4 & Aceptar \\ \hline
	5 a 15 & Mitigar \\ \hline
	16 a 25 & Evitar \\ \hline
	
\end{tabular}
\end{center}

\section*{Presupuesto}

Duración total del proyecto: 480hs.

\begin{itemize}
\item Estación de trabajo: \$16000.
\item Servidor de desarrollo: \$10000.
\item Conexión a internet: \$3600.
\item Horas de trabajo (\$40 por hora): \$19200.
\end{itemize}

Costo total: \$48800.


\newpage
 
\nocite{*}
\printbibliography
\end{document}




